\begin{Exercise}[label={ex:indici}]
  Scrivere un programma che acquisisca un certo numero di valori stabilito dall'utente (massimo $10$), e poi acquisisca un eguale quantit\`a di indici, che specificano l'ordine in cui il programma dovr\`a stampare i valori acquisiti in precedenza. Per esempio:

\begin{lstlisting}[style=console]
Quanti elementi vuoi inserire (max: 10)? 5
Inserire il 1o valore: 10
Inserire il 2o valore: 20
Inserire il 3o valore: 30
Inserire il 4o valore: 40
Inserire il 5o valore: 50
Qual è il 1o valore che vuoi stampare? 3
Qual è il 2o valore che vuoi stampare? 2
Qual è il 3o valore che vuoi stampare? 1
Qual è il 4o valore che vuoi stampare? 4
Qual è il 5o valore che vuoi stampare? 5
30
20
10
40
50
\end{lstlisting}
\end{Exercise}

\begin{Answer}[ref={ex:indici}]
  \lstinputlisting[style=cstylesmall]{src/indice.c}
\end{Answer}
